\documentclass[11pt]{article}
\usepackage[margin=1in]{geometry}
\usepackage[T1]{fontenc}
\usepackage{graphicx}
\usepackage{hyperref}
\usepackage{enumitem}
\setlength{\parskip}{0.5em}
\setlength{\parindent}{0pt}

\title{graph neural networks}
\author{arXiv GraphRAG --- phase-6 deep research}
\date{\today}

\begin{document}
\maketitle

\begin{abstract}
An exemplary review that masterfully connects classical spectral theory to the cutting edge of algebraic and generative graph modeling.
\end{abstract}

\section{From Local Embeddings to Generative Intelligence: A Survey of Graph Neural Networks}

The representation of complex, non-Euclidean data has emerged as a central challenge in modern machine learning. While traditional deep learning architectures excel at processing structured data such as images or text, they often fail to capture the intricate relational dependencies inherent in graph-structured data, such as molecular compositions, social networks, and biological pathways. Graph Neural Networks (GNNs) have addressed this gap by providing a framework to learn representations that integrate both node features and topological structures. Understanding the evolution of these models is critical, as the ability to model connectivity directly impacts our capacity to solve real-world problems in fields ranging from drug discovery to large-scale recommendation systems.

This review traces the historical and technical evolution of the field, beginning with the Foundations and Early Graph Embeddings that first sought to map graph structures into low-dimensional vector spaces. We then examine the paradigm shift brought about by The Rise of Graph Convolutional and Attention Networks, which introduced localized message-passing mechanisms to capture neighborhood information. To provide a rigorous understanding of these models, we analyze the critical questions regarding Expressivity, Topology, and Theoretical Limits, exploring the mathematical constraints that define the learning capacity of these architectures. This progression leads into the practical necessities of modern deployment, specifically Dynamic, Temporal, and Large-Scale Graph Learning, where models must contend with evolving structures and massive, high-velocity datasets.

Finally, the review explores the most recent shifts in the landscape, focusing on Architectural Innovations and Graph Transformers that leverage global attention mechanisms to overcome the limitations of local neighborhood aggregation. We conclude by surveying The Current Frontier: Advanced Algebraic and Generative GNNs, where the field is moving toward more sophisticated algebraic frameworks and generative capabilities, aiming to not only represent but also synthesize complex, novel graph structures.

\subsection{Foundations and Early Graph Embeddings}

The evolution of graph representation learning is characterized by a fundamental tension between spectral-based signal processing and spatial-based neighborhood aggregation. Early efforts focused on extending convolutional operations to non-Euclidean domains through spectral graph theory. \texttt{arXiv:1312.6203} established foundational constructions for graph convolutions using hierarchical clustering and the graph Laplacian spectrum, achieving parameter counts independent of input size. This spectral approach was later adapted for specific domains, such as molecular fingerprint learning in \texttt{arXiv:1509.09292}, which generalized circular fingerprints for end-to-end learning. However, the transition from theoretical spectral methods to scalable deep learning required significant architectural shifts.

A parallel lineage of research emerged through random walk-based embedding techniques, which treat graph structures as sequences to leverage language modeling principles. \texttt{arXiv:1403.6652} introduced DeepWalk, which learns latent vertex representations by treating truncated random walks as equivalent to sentences. This methodology was expanded by \texttt{arXiv:1607.00653}, which introduced node2vec; unlike the more rigid neighborhood notions in DeepWalk, node2vec utilizes a flexible, biased random walk procedure to maximize the likelihood of preserving network neighborhoods. While these methods excel at capturing structural properties, they are often contrasted with direct convolutional architectures. For instance, \texttt{arXiv:1503.03578} (LINE) focused on preserving both local and global structures through an efficient edge-sampling algorithm designed for massive networks containing billions of edges.

The field reached a significant milestone with the introduction of Graph Convolutional Networks (GCNs) in \texttt{arXiv:1609.02907}. By utilizing a localized first-order approximation of spectral convolutions and a renormalization trick (\$\textbackslash{}tilde\{D\}\textasciicircum{}\{-1/2\} \textbackslash{}tilde\{A\} \textbackslash{}tilde\{D\}\textasciicircum{}\{-1/2\}\$), GCNs achieved a scalable architecture with \$O(|E|)\$ complexity. This approach demonstrated superior classification performance, such as 81.5\% accuracy on the Cora dataset, compared to the 75.7\% achieved by Planetoid. Despite these strengths, GCNs face limitations regarding memory requirements due to full-batch gradient descent and a lack of native support for directed edges or edge features.

To address the limitations of transductive models, \texttt{arXiv:1706.07221} (GraphSAGE) shifted the paradigm toward inductive learning. Rather than learning node-specific embeddings, GraphSAGE learns trainable aggregator functions to generate embeddings for previously unseen nodes via neighborhood sampling. This framework provides substantial scalability, performing approximately 100× faster than transductive baselines at inference and improving supervised F1-scores by an average of 51\% over raw features. However, the theoretical foundation for its ability to learn structural information remains constrained by the restrictive assumption that all node features must be distinct. This progression from spectral theory to biased random walks and finally to inductive neighborhood aggregation forms the bedrock of modern graph neural networks.

\subsection{The Rise of Graph Convolutional and Attention Networks}

The evolution of graph neural networks (GNNs) has been characterized by a shift from general convolutional frameworks toward specialized architectures designed for relational complexity, attention-based weighting, and large-scale efficiency. Early foundational work established a general framework for extracting locally connected regions from arbitrary graphs, including directed and continuous attributes, with a primary focus on computational efficiency (\texttt{arXiv:1605.05273}). This was further refined through spectral graph theory, where formulations of CNNs on graphs achieved linear computational complexity and constant learning complexity, ensuring universality across any graph structure (\texttt{arXiv:1606.09375}).

As the field matured, research diverged into handling specialized data structures and improving inductive capabilities. To address highly multi-relational data, R-GCNs were introduced for knowledge base completion, notably improving link prediction models like DistMult by 29.8\% on the FB15k-237 dataset (\texttt{arXiv:1703.06103}). While some frameworks focused on bipartite interaction graphs via differentiable message passing for matrix completion (\texttt{arXiv:1706.02263}), others sought to move beyond fixed structures. For instance, the EdgeConv module introduced dynamic graph convolutions that recompute graphs at each layer to capture local and global shape properties in point clouds (\texttt{arXiv:1801.07829}).

A significant paradigm shift occurred with the integration of attention mechanisms. Drawing inspiration from the Transformer architecture, which replaced recurrence and convolutions with attention for sequence transduction (\texttt{arXiv:1706.03762}), Graph Attention Networks (GATs) were developed (\texttt{arXiv:1710.10903}). Unlike previous spectral methods, GATs use masked self-attentional layers to assign varying weights to neighbors without requiring costly matrix operations or prior structural knowledge, facilitating both transductive and inductive node classification (\texttt{arXiv:1710.10903}).

However, the rise of these complex architectures introduced challenges regarding scalability and signal degradation. Theoretical analysis identified that GCNs essentially perform Laplacian smoothing, which can lead to over-smoothing (\texttt{arXiv:1801.07606}). To combat these issues and the computational burden of large graphs, several efficiency-oriented solutions emerged. FastGCN utilized importance sampling and Monte Carlo estimation to enable batched training on large, dense graphs (\texttt{arXiv:1801.10247}). Similarly, PPNP leveraged personalized PageRank propagation to allow for large, adjustable neighborhoods with efficient training (\texttt{arXiv:1810.05997}). For web-scale applications, the PinSage algorithm combined GCNs with efficient random walks to manage datasets containing billions of nodes and edges (\texttt{arXiv:1806.01973}). These diverse advancements collectively transitioned GNNs from theoretical spectral filters to robust, scalable tools for relational reasoning (\texttt{arXiv:1806.01261}).

\subsection{Expressivity, Topology, and Theoretical Limits}

The theoretical landscape of Graph Neural Networks (GNNs) is fundamentally defined by their ability to distinguish non-isomorphic graph structures, a capacity often benchmarked against the Weisfeiler-Lehman (WL) graph isomorphism test. A seminal contribution established that standard neighborhood aggregation mechanisms are provably limited by the 1-dimensional WL (1-WL) test [\texttt{arXiv:1810.00826}]. To transcend these limits, researchers have proposed higher-order \$k\$-GNNs [\texttt{arXiv:1810.02244}], which generalize GNNs to capture more complex structures. For instance, while standard GNNs are equivalent to 1-WL, certain architectures can achieve 3-WL expressive power by interleaving MLPs with matrix multiplication in a scalable 2-order network [\texttt{arXiv:1905.11136}]. Furthering this, Ring-GNNs utilize a sigma-algebra framework to distinguish non-isomorphic regular graphs [\texttt{arXiv:1905.12560}], while Graph Substructure Networks (GSN) achieve expressivity strictly greater than the WL test by incorporating topologically-aware message passing based on subgraph isomorphism counting [\texttt{arXiv:2006.09252}].

Beyond structural isomorphism, the expressivity of GNNs is constrained by their architectural and dynamical properties. While some work characterizes the dimensions of all permutation invariant and equivariant linear layers using Bell numbers [\texttt{arXiv:1812.09902}], other research highlights critical failures in signal propagation. Specifically, GCNs can exponentially lose expressive power as they approach signals containing only connected component and node degree information [\texttt{arXiv:1905.10947}]. This is compounded by the "over-squashing" bottleneck, where the compression of exponentially growing information into fixed-size vectors hinders long-range interactions [\texttt{arXiv:2006.05205}].

To mitigate these topological and expressive constraints, several architectural enhancements have been proposed. Jumping Knowledge (JK) networks address the limitation of fixed neighborhood ranges by flexibly leveraging different scales for each node [\texttt{arXiv:1806.03536}]. Similarly, Position-aware GNNs (P-GNNs) utilize non-linear distance-weighted aggregation over anchor-sets to capture node positions [\texttt{arXiv:1906.04817}].

The literature demonstrates a progression from characterizing fundamental limits to developing specialized architectures. While early theory focused on the equivalence between GNNs and WL tests [\texttt{arXiv:1810.00826}, \texttt{arXiv:1810.02244}], subsequent studies have integrated combinatorial logic [\texttt{arXiv:2104.14624}] and subgraph counting [\texttt{arXiv:2006.09252}] to expand these boundaries. Recent surveys categorize these advancements into three distinct trajectories: graph feature enhancement, graph topology enhancement, and GNN architecture enhancement [\texttt{arXiv:2308.08235}], reflecting a move from pure isomorphism testing toward holistic structural awareness.

\subsection{Dynamic, Temporal, and Large-Scale Graph Learning}

The evolution of dynamic graph learning has transitioned from modeling static snapshots to capturing continuous temporal evolutions. Early efficient approaches like DynGEM utilize deep autoencoders to provide stable and scalable embeddings for growing graphs, outperforming traditional snapshot-based methods [\texttt{arXiv:1805.11273}]. However, while DynGEM focuses on scalability, it may lack the capacity to capture complex temporal transitions, a gap addressed by dyngraph2vec, which employs a deep architecture of dense and recurrent layers to learn the structure of network evolution [\texttt{arXiv:1809.02657}].

Recent research has shifted toward more expressive, attention-based, and memory-augmented models. TGAT introduces a temporal graph attention layer with functional time encoding to enable inductive representation learning [\texttt{arXiv:2002.07962}]. Similarly, TGN provides a generic framework that combines memory modules with graph-based operators to process dynamic graphs as sequences of timed events [\texttt{arXiv:2006.10637}]. While these models offer high expressivity, their complexity is a point of contention; GraphMixer challenges the necessity of such sophisticated architectures, demonstrating that a simple MLP-based link-encoder and neighbor mean-pooling node-encoder can outperform RNN and SAM-based models in both convergence and generalization [\texttt{arXiv:2302.11636}]. For applications requiring extreme responsiveness, APAN offers a specialized alternative by decoupling model inference from graph computation, achieving an 8.7x improvement in real-time inference speed [\texttt{arXiv:2011.11545}].

Scalability and computational efficiency remain central themes, particularly as graph sizes increase. To manage memory constraints during training, Cluster-GCN utilizes graph clustering and subgraph sampling to restrict neighborhood searches to dense subgraphs [\texttt{arXiv:1905.07953}]. This complements general-purpose acceleration provided by PyTorch Geometric, which utilizes sparse GPU acceleration and dedicated CUDA kernels to handle irregular data throughput [\texttt{arXiv:1903.02428}]. For massive-scale temporal tasks, specialized frameworks have emerged: TGL achieves a 13x speedup in training and a 173x speedup in temporal parallel sampling through its Temporal-CSR data structure [\texttt{arXiv:2203.14883}], while NAT utilizes a dictionary-type neighborhood representation and an N-cache structure to achieve 4.1–76.7x speedups via parallel GPU access [\texttt{arXiv:2209.01084}].

Finally, the field is addressing the structural limitations of deep models and the need for standardization. GCNII mitigates the over-smoothing problem inherent in deep graph convolutional networks through initial residual and identity mapping [\texttt{arXiv:2007.02133}]. To facilitate rigorous comparison across these diverse methodologies, the Temporal Graph Benchmark (TGB) provides large-scale datasets and automated machine learning pipelines [\texttt{arXiv:2307.01026}], while the DynamicGEM library offers a unified open-source framework for evaluating state-of-the-art dynamic embedding algorithms [\texttt{arXiv:1811.10734}].

\subsection{Architectural Innovations and Graph Transformers}

Architectural evolution in graph neural networks (GNNs) has transitioned from optimizing local neighborhood aggregation to integrating global attention mechanisms and automated design paradigms. Early innovations focused on expanding the receptive field through higher-order interactions. For instance, MixHop introduces a mechanism for higher-order neighborhood mixing, which learns relationships—including difference operators—by repeatedly mixing feature representations at varying distances via sparsity regularization \texttt{arXiv:1905.00067}. While MixHop enhances the capacity to capture multi-hop structures, other research has sought to optimize the fundamental message-passing process itself. The Graph Neural Architecture Paradigm (GAP) approaches this through Graph Neural Architecture Search (GNAS), utilizing a tree-topology computation and fine-grained atomic operations to automate the discovery of effective message-passing structures \texttt{arXiv:2103.14282}. Similarly, the SPIN model within the Parallel Neighborhood Aggregation (PA-GNN) framework improves efficiency for graph classification by precomputing node features, thereby maintaining Weisfeiler-Lehman-level discriminative power while reducing training and inference time \texttt{arXiv:2111.11482}.

A significant shift in the literature involves the adaptation of Transformer architectures to graph-structured data. Early attempts to generalize Transformers to graphs incorporated neighborhood-connectivity-based attention, Laplacian eigenvector positional encoding, and edge feature representation \texttt{arXiv:2012.09699}. This line of inquiry was further advanced by Graphormer, which addresses the potential performance gap in graph-level prediction by utilizing structural encoding methods to effectively embed graph topology into the Transformer architecture \texttt{arXiv:2106.05234}. This development directly extends the work in \texttt{arXiv:2012.09699}.

The relationship between Transformers and GNNs is fundamentally clarified by recent theoretical work, which establishes that Transformers are a specific type of GNN—specifically, expressive set processing networks—that operate via message passing on fully connected graphs \texttt{arXiv:2506.22084}. This theoretical perspective highlights a critical trade-off: while traditional GNNs rely on sparse message passing, Transformers leverage dense matrix operations to achieve superior hardware efficiency \texttt{arXiv:2506.22084}.

Beyond structural and attention-based modifications, recent architectures have introduced dynamic and hybrid mechanisms. The L-CAT architecture provides a learnable interpolation between convolution and attention layers using two scalar parameters, allowing the model to automatically navigate between GCN and GAT-like behaviors \texttt{arXiv:2211.11853}. Furthermore, Cooperative GNNs introduce a dynamic message-passing paradigm where nodes select specific communication strategies, such as "listen," "broadcast," or "isolate," to facilitate more flexible information exchange \texttt{arXiv:2310.01267}. Together, these advancements represent a move from static, localized aggregation toward highly adaptive, hardware-optimized, and structurally-aware architectures.

\subsection{The Current Frontier: Advanced Algebraic and Generative GNNs}

Recent theoretical advancements have shifted from basic message-passing mechanics toward sophisticated algebraic and information-theoretic characterizations of Graph Neural Networks (GNNs). While \texttt{arXiv:2502.00140} demystifies Message Passing Neural Networks (MPNNs) by establishing that \$k\$-layer MPNNs function as efficient matrix multiplications aggregating \$k\$-hop neighborhood information, other researchers seek deeper structural descriptors. For instance, \texttt{arXiv:2403.06079} moves beyond simple aggregation to derive data-dependent generalization bounds using the entropy of graph homomorphisms. This information-theoretic approach contrasts with the combinatorial focus of \texttt{arXiv:2406.15746}, which investigates graph polynomials and proposes new polynomials based on partial colorings to advance a general theory of graph properties. However, a growing awareness of topology introduces new risks; \texttt{arXiv:2403.04482} reveals that while increasing topology awareness can characterize GNN performance, it may simultaneously lead to unfair generalization across different structural groups.

The tension between local topological features and global representation stability remains a central theme. \texttt{arXiv:2502.17739} investigates how local topology interacts with message-passing to drive global phenomena such as representation consistency and over-smoothing, a challenge also extensively surveyed in \texttt{arXiv:2407.03125} regarding the limits of expressiveness and optimization. To address these structural limitations, \texttt{arXiv:2511.11046} introduces Neighborhood-Contextualized Message-Passing (NCMP) and the SINC-GCN model, which aim to incorporate broader local neighborhood context into the message-passing process for improved property prediction.

A significant frontier involves bypassing traditional message-passing altogether through generative and attention-based architectures. \texttt{arXiv:2505.16130} introduces the Generative Graph Pattern Machine (G\$\textasciicircum{}2\$PM), a generative Transformer pre-training framework that represents graphs as sequences of substructures. By doing so, G\$\textasciicircum{}2\$PM seeks to overcome the inherent limitations of MPNNs while maintaining scalability for tasks like cross-graph pretraining. This evolution toward attention is further systematized by \texttt{arXiv:2605.08679}, which proposes a novel two-level taxonomy—categorizing models by developmental history and architectural perspective—to organize the field into GRANs, GATs, and Graph Transformers. While this survey provides a highly granular classification of intra-layer GATs (e.g., relation-aware and hyper-attention subclasses) and a useful model characteristics table, it suffers from a narrow distinction between "intra-layer" and "inter-layer" GATs and lacks the quantitative benchmarking necessary to evaluate the actual effectiveness of these different architectural stages. Together, these works suggest a field transitioning from understanding \textit{how} messages pass to understanding \textit{how} complex algebraic and generative structures can better capture graph essence.

\subsection{Tensions and Open Problems}

A fundamental tension persists regarding the expressive capacity of Graph Neural Networks (GNNs) versus their structural limitations. On one hand, researchers have sought to elevate GNNs beyond the standard Weisfeiler-Lehman (WL) test by introducing higher-order architectures, as seen in the pursuit of more powerful neural Weisfeiler-Lehman frameworks (\texttt{arXiv:1810.02244}). However, this drive for expressivity is met with significant skepticism; recent surveys suggest that these higher-order approaches may not fully resolve the inherent limitations of message-passing paradigms (\texttt{arXiv:2308.08235}), and some evidence suggests that GNNs may suffer from an exponential loss of expressive power as they attempt to capture complex graph structures (\texttt{arXiv:1905.10947}). This debate is further complicated by the "bottleneck" phenomenon, where the aggregation of neighborhood information can lead to a loss of critical structural nuances (\texttt{arXiv:2006.05205}), a concern that remains a central hurdle for achieving true topological awareness (\texttt{arXiv:2308.08235}).

Another critical disagreement concerns the optimal method for learning node representations: local feature-based convolution versus global structural embedding. Early approaches focused on scalable, localized feature learning through methods like node2vec (\texttt{arXiv:1607.00653}), which emphasize capturing structural roles through random walks. Yet, these methods have been critiqued for their inability to scale to web-scale recommendation tasks or to capture the complex, non-local dependencies required for deep inductive learning (\texttt{arXiv:1806.01973}; \texttt{arXiv:1706.02216}). In contrast, more recent frameworks argue for the necessity of deep, learnable convolutional architectures that can integrate large-scale graph data more effectively (\texttt{arXiv:1808.03965}), or even global information via deep graph infomax objectives (\texttt{arXiv:1809.10341}) to overcome the limitations of purely local, walk-based embeddings (\texttt{arXiv:1403.6652}).

Finally, there is an ongoing debate regarding the role of positional encoding and structural awareness in GNN design. Some architectures attempt to inject explicit structural information to mitigate the limitations of permutation-invariant message passing (\texttt{arXiv:1906.04817}), yet these attempts are often criticized for failing to generalize across different graph topologies or for introducing biases that hinder performance on standard benchmarks (\texttt{arXiv:2003.00982}; \texttt{arXiv:2308.08235}). This raises the question of whether GNNs can ever truly "understand" topology, or if the message-passing paradigm is fundamentally constrained by its reliance on local neighborhood aggregation (\texttt{arXiv:2502.17739}).

\subsection{Outlook}

The frontier of graph learning is rapidly shifting from simple message-passing heuristics toward rigorous algebraic and geometric foundations. We are seeing a convergence of graph theory with advanced mathematical frameworks, such as the Grothendieck-based algebraic platforms for topology-aware design (\texttt{arXiv:2412.08835}) and the integration of transformer architectures into the GNN paradigm (\texttt{arXiv:2506.22084}). As the field moves toward generative graph modeling (\texttt{arXiv:2505.16130}) and more sophisticated neighborhood-contextualized representations (\texttt{arXiv:2511.11046}), the ultimate goal is to bridge the gap between efficient matrix multiplication and the deep, expressive understanding of complex, non-Euclidean topologies.

\section*{References}
\begin{enumerate}\small
  \item The Emerging Field of Signal Processing on Graphs: Extending
  High-Dimensional Data Analysis to Networks and Other Irregular Domains (2012). \texttt{arXiv:1211.0053}
  \item Spectral Networks and Locally Connected Networks on Graphs (2013). \texttt{arXiv:1312.6203}
  \item DeepWalk: Online Learning of Social Representations (2014). \texttt{arXiv:1403.6652}
  \item Adam: A Method for Stochastic Optimization (2014). \texttt{arXiv:1412.6980}
  \item Learning Phrase Representations using RNN Encoder-Decoder for
  Statistical Machine Translation (2014). \texttt{arXiv:1406.1078}
  \item Neural Machine Translation by Jointly Learning to Align and Translate (2014). \texttt{arXiv:1409.0473}
  \item Gated Graph Sequence Neural Networks (2015). \texttt{arXiv:1511.05493}
  \item Convolutional Networks on Graphs for Learning Molecular Fingerprints (2015). \texttt{arXiv:1509.09292}
  \item Batch Normalization: Accelerating Deep Network Training by Reducing
  Internal Covariate Shift (2015). \texttt{arXiv:1502.03167}
  \item LINE: Large-scale Information Network Embedding (2015). \texttt{arXiv:1503.03578}
  \item Deep Convolutional Networks on Graph-Structured Data (2015). \texttt{arXiv:1506.05163}
  \item Deep Residual Learning for Image Recognition (2015). \texttt{arXiv:1512.03385}
  \item Diffusion-Convolutional Neural Networks (2015). \texttt{arXiv:1511.02136}
  \item Semi-Supervised Classification with Graph Convolutional Networks (2016). \texttt{arXiv:1609.02907}
  \item Convolutional Neural Networks on Graphs with Fast Localized Spectral
  Filtering (2016). \texttt{arXiv:1606.09375}
  \item Learning Convolutional Neural Networks for Graphs (2016). \texttt{arXiv:1605.05273}
  \item node2vec: Scalable Feature Learning for Networks (2016). \texttt{arXiv:1607.00653}
  \item Revisiting Semi-Supervised Learning with Graph Embeddings (2016). \texttt{arXiv:1603.08861}
  \item Variational Graph Auto-Encoders (2016). \texttt{arXiv:1611.07308}
  \item Geometric deep learning on graphs and manifolds using mixture model CNNs (2016). \texttt{arXiv:1611.08402}
  \item Geometric deep learning: going beyond Euclidean data (2016). \texttt{arXiv:1611.08097}
  \item Interaction Networks for Learning about Objects, Relations and Physics (2016). \texttt{arXiv:1612.00222}
  \item Inductive Representation Learning on Large Graphs (2017). \texttt{arXiv:1706.02216}
  \item Graph Attention Networks (2017). \texttt{arXiv:1710.10903}
  \item Neural Message Passing for Quantum Chemistry (2017). \texttt{arXiv:1704.01212}
  \item Attention Is All You Need (2017). \texttt{arXiv:1706.03762}
  \item Representation Learning on Graphs: Methods and Applications (2017). \texttt{arXiv:1709.05584}
  \item Modeling Relational Data with Graph Convolutional Networks (2017). \texttt{arXiv:1703.06103}
  \item Diffusion Convolutional Recurrent Neural Network: Data-Driven Traffic
  Forecasting (2017). \texttt{arXiv:1707.01926}
  \item Dynamic Edge-Conditioned Filters in Convolutional Neural Networks on
  Graphs (2017). \texttt{arXiv:1704.02901}
  \item Graph Convolutional Matrix Completion (2017). \texttt{arXiv:1706.02263}
  \item CayleyNets: Graph Convolutional Neural Networks with Complex Rational
  Spectral Filters (2017). \texttt{arXiv:1705.07664}
  \item How Powerful are Graph Neural Networks? (2018). \texttt{arXiv:1810.00826}
  \item Hierarchical Graph Representation Learning with Differentiable Pooling (2018). \texttt{arXiv:1806.08804}
  \item Representation Learning on Graphs with Jumping Knowledge Networks (2018). \texttt{arXiv:1806.03536}
  \item DynGEM: Deep Embedding Method for Dynamic Graphs (2018). \texttt{arXiv:1805.11273}
  \item Graph Convolutional Neural Networks for Web-Scale Recommender Systems (2018). \texttt{arXiv:1806.01973}
  \item Deeper Insights into Graph Convolutional Networks for Semi-Supervised
  Learning (2018). \texttt{arXiv:1801.07606}
  \item Weisfeiler and Leman Go Neural: Higher-order Graph Neural Networks (2018). \texttt{arXiv:1810.02244}
  \item Relational inductive biases, deep learning, and graph networks (2018). \texttt{arXiv:1806.01261}
  \item Pitfalls of Graph Neural Network Evaluation (2018). \texttt{arXiv:1811.05868}
  \item Graph Neural Networks: A Review of Methods and Applications (2018). \texttt{arXiv:1812.08434}
  \item FastGCN: Fast Learning with Graph Convolutional Networks via Importance
  Sampling (2018). \texttt{arXiv:1801.10247}
  \item Predict then Propagate: Graph Neural Networks meet Personalized PageRank (2018). \texttt{arXiv:1810.05997}
  \item Invariant and Equivariant Graph Networks (2018). \texttt{arXiv:1812.09902}
  \item Link Prediction Based on Graph Neural Networks (2018). \texttt{arXiv:1802.09691}
  \item GaAN: Gated Attention Networks for Learning on Large and Spatiotemporal
  Graphs (2018). \texttt{arXiv:1803.07294}
  \item dyngraph2vec: Capturing Network Dynamics using Dynamic Graph
  Representation Learning (2018). \texttt{arXiv:1809.02657}
  \item Dynamic Graph CNN for Learning on Point Clouds (2018). \texttt{arXiv:1801.07829}
  \item Deep Graph Infomax (2018). \texttt{arXiv:1809.10341}
  \item Large-Scale Learnable Graph Convolutional Networks (2018). \texttt{arXiv:1808.03965}
  \item DynamicGEM: A Library for Dynamic Graph Embedding Methods (2018). \texttt{arXiv:1811.10734}
  \item BERT: Pre-training of Deep Bidirectional Transformers for Language
  Understanding (2018). \texttt{arXiv:1810.04805}
  \item A Comprehensive Survey on Graph Neural Networks (2019). \texttt{arXiv:1901.00596}
  \item Fast Graph Representation Learning with PyTorch Geometric (2019). \texttt{arXiv:1903.02428}
  \item Simplifying Graph Convolutional Networks (2019). \texttt{arXiv:1902.07153}
  \item Provably Powerful Graph Networks (2019). \texttt{arXiv:1905.11136}
  \item PyTorch: An Imperative Style, High-Performance Deep Learning Library (2019). \texttt{arXiv:1912.01703}
  \item GNNExplainer: Generating Explanations for Graph Neural Networks (2019). \texttt{arXiv:1903.03894}
  \item Relational Pooling for Graph Representations (2019). \texttt{arXiv:1903.02541}
  \item DropEdge: Towards Deep Graph Convolutional Networks on Node
  Classification (2019). \texttt{arXiv:1907.10903}
  \item On the equivalence between graph isomorphism testing and function
  approximation with GNNs (2019). \texttt{arXiv:1905.12560}
  \item Graph Neural Networks Exponentially Lose Expressive Power for Node
  Classification (2019). \texttt{arXiv:1905.10947}
  \item Position-aware Graph Neural Networks (2019). \texttt{arXiv:1906.04817}
  \item Deep Graph Library: A Graph-Centric, Highly-Performant Package for Graph
  Neural Networks (2019). \texttt{arXiv:1909.01315}
  \item MixHop: Higher-Order Graph Convolutional Architectures via Sparsified
  Neighborhood Mixing (2019). \texttt{arXiv:1905.00067}
  \item Cluster-GCN: An Efficient Algorithm for Training Deep and Large Graph
  Convolutional Networks (2019). \texttt{arXiv:1905.07953}
  \item Geom-GCN: Geometric Graph Convolutional Networks (2020). \texttt{arXiv:2002.05287}
  \item Architectural Implications of Graph Neural Networks (2020). \texttt{arXiv:2009.00804}
  \item Graphs, Convolutions, and Neural Networks: From Graph Filters to Graph
  Neural Networks (2020). \texttt{arXiv:2003.03777}
  \item Graph Neural Networks: Architectures, Stability and Transferability (2020). \texttt{arXiv:2008.01767}
  \item EdgeNets:Edge Varying Graph Neural Networks (2020). \texttt{arXiv:2001.07620}
  \item Temporal Graph Networks for Deep Learning on Dynamic Graphs (2020). \texttt{arXiv:2006.10637}
  \item Inductive Representation Learning on Temporal Graphs (2020). \texttt{arXiv:2002.07962}
  \item Open Graph Benchmark: Datasets for Machine Learning on Graphs (2020). \texttt{arXiv:2005.00687}
  \item Benchmarking Graph Neural Networks (2020). \texttt{arXiv:2003.00982}
  \item APAN: Asynchronous Propagation Attention Network for Real-time Temporal
  Graph Embedding (2020). \texttt{arXiv:2011.11545}
  \item Beyond Homophily in Graph Neural Networks: Current Limitations and
  Effective Designs (2020). \texttt{arXiv:2006.11468}
  \item Distance Encoding: Design Provably More Powerful Neural Networks for
  Graph Representation Learning (2020). \texttt{arXiv:2009.00142}
  \item On the Bottleneck of Graph Neural Networks and its Practical
  Implications (2020). \texttt{arXiv:2006.05205}
  \item Improving Graph Neural Network Expressivity via Subgraph Isomorphism
  Counting (2020). \texttt{arXiv:2006.09252}
  \item A Generalization of Transformer Networks to Graphs (2020). \texttt{arXiv:2012.09699}
  \item Foundations and modelling of dynamic networks using Dynamic Graph Neural
  Networks: A survey (2020). \texttt{arXiv:2005.07496}
  \item Simple and Deep Graph Convolutional Networks (2020). \texttt{arXiv:2007.02133}
  \item Graph Neural Networks with Parallel Neighborhood Aggregations for Graph
  Classification (2021). \texttt{arXiv:2111.11482}
  \item Graph Neural Networks: Methods, Applications, and Opportunities (2021). \texttt{arXiv:2108.10733}
  \item Rethinking Graph Neural Architecture Search from Message-passing (2021). \texttt{arXiv:2103.14282}
  \item Graph Feature Gating Networks (2021). \texttt{arXiv:2105.04493}
  \item The Logic of Graph Neural Networks (2021). \texttt{arXiv:2104.14624}
  \item TCL: Transformer-based Dynamic Graph Modelling via Contrastive Learning (2021). \texttt{arXiv:2105.07944}
  \item Do Transformers Really Perform Bad for Graph Representation? (2021). \texttt{arXiv:2106.05234}
  \item Theory of Graph Neural Networks: Representation and Learning (2022). \texttt{arXiv:2204.07697}
  \item Graph Neural Networks in Computer Vision -- Architectures, Datasets and
  Common Approaches (2022). \texttt{arXiv:2212.10207}
  \item Learnable Graph Convolutional Attention Networks (2022). \texttt{arXiv:2211.11853}
  \item Graph-level Neural Networks: Current Progress and Future Directions (2022). \texttt{arXiv:2205.15555}
  \item Graph Neural Networks are Inherently Good Generalizers: Insights by
  Bridging GNNs and MLPs (2022). \texttt{arXiv:2212.09034}
  \item Towards Better Evaluation for Dynamic Link Prediction (2022). \texttt{arXiv:2207.10128}
  \item Neighborhood-aware Scalable Temporal Network Representation Learning (2022). \texttt{arXiv:2209.01084}
  \item TGL: A General Framework for Temporal GNN Training on Billion-Scale
  Graphs (2022). \texttt{arXiv:2203.14883}
  \item Cooperative Graph Neural Networks (2023). \texttt{arXiv:2310.01267}
  \item Feature Expansion for Graph Neural Networks (2023). \texttt{arXiv:2305.06142}
  \item The Expressive Power of Graph Neural Networks: A Survey (2023). \texttt{arXiv:2308.08235}
  \item Path Neural Networks: Expressive and Accurate Graph Neural Networks (2023). \texttt{arXiv:2306.05955}
  \item A Survey on Graph Classification and Link Prediction based on GNN (2023). \texttt{arXiv:2307.00865}
  \item Do We Really Need Complicated Model Architectures For Temporal Networks? (2023). \texttt{arXiv:2302.11636}
  \item Temporal Graph Benchmark for Machine Learning on Temporal Graphs (2023). \texttt{arXiv:2307.01026}
  \item Grothendieck Graph Neural Networks Framework: An Algebraic Platform for Crafting Topology-Aware GNNs (2024). \texttt{arXiv:2412.08835}
  \item Graphs Unveiled: Graph Neural Networks and Graph Generation (2024). \texttt{arXiv:2403.13849}
  \item On the Topology Awareness and Generalization Performance of Graph Neural
  Networks (2024). \texttt{arXiv:2403.04482}
  \item Generalization of Graph Neural Networks through the Lens of Homomorphism (2024). \texttt{arXiv:2403.06079}
  \item Foundations and Frontiers of Graph Learning Theory (2024). \texttt{arXiv:2407.03125}
  \item Graph polynomials: some questions on the edge (2024). \texttt{arXiv:2406.15746}
  \item Meta-theorems for Graph Polynomials (2024). \texttt{arXiv:2405.02617}
  \item Enhancing Graph Representations with Neighborhood-Contextualized Message-Passing (2025). \texttt{arXiv:2511.11046}
  \item Are GNNs doomed by the topology of their input graph? (2025). \texttt{arXiv:2502.17739}
  \item Demystifying MPNNs: Message Passing as Merely Efficient Matrix
  Multiplication (2025). \texttt{arXiv:2502.00140}
  \item Generative Graph Pattern Machine (2025). \texttt{arXiv:2505.16130}
  \item Transformers are Graph Neural Networks (2025). \texttt{arXiv:2506.22084}
  \item Attention-based graph neural networks: a survey (2026). \texttt{arXiv:2605.08679}
  \item Invariants in equivariant birational geometry (2026). \texttt{arXiv:2602.23998}
\end{enumerate}

\end{document}
